\subsection{扩张叠}\label{src:9.2.4-subsection}
我们讨论参数化两个给定向量丛的扩张的模函子与模叠的代数性. 由此可在 $\cBun_{r_1,d_1}\times\cBun_{r_2,d_2}$ 上构造扩张叠, 它有到 $\cBun_{r_1+r_2,d_1+d_2}$ 的态射. 这可视为类似于粘合态射 (\bookref{6.6.16}) $\mathcal M_{g_1,n_1}\times\mathcal M_{g_2,n_2}\to\mathcal M_{g_1+g_2,n_1+n_2-2}$ 的构造. 与稳定曲线的情形一样, 这提供了建立 $\cBun_{r,d}$ 的性质的归纳策略. 下一小节将利用这一策略证明稳定轨迹 $\cBun_{r,d}^{\mathrm s}$ 非空.

\textbf{扩张.} 若 $F_1,F_2$ 是概形 $X$ 上的向量丛, 以 $\mathrm{EXT}(F_1,F_2)$ 表示扩张 $0\to F_2\xrightarrow aE\xrightarrow bF_1\to0$ 的群胚, 其中扩张之间的态射是嵌入如下交换图的同构 $\alpha:E\to E$:
\[
\begin{tikzcd}[ampersand replacement=\&,column sep=small]
0\arrow[r]\&F_2\arrow[r,"a"]\arrow[d,"\mathrm{id}"']\&E\arrow[r,"b"]\arrow[d,"\alpha"]\&F_1\arrow[r]\arrow[d,"\mathrm{id}"]\&0\\
0\arrow[r]\&F_2\arrow[r,"a"]\&E'\arrow[r,"b"]\&F_1\arrow[r]\&0.
\end{tikzcd}
\]
\begin{exercise}[简单]\label{src:9.2.10}
证明 $\operatorname{Aut}_{\mathrm{EXT}(F_1,F_2)}([E])\cong\Hom_{\mathcal O_X}(F_1,F_2)$, 其中 $\gamma\in\Hom_{\mathcal O_X}(F_1,F_2)$ 对应于扩张 $E$ 的自同构 $\mathrm{id}+a\circ\gamma\circ b$.
\end{exercise}
\textbf{扩张函子.} 讨论扩张叠之前, 先描述扩张函子的性质. 回顾群 $\operatorname{Ext}_{\mathcal O_X}^1(F_1,F_2)$ 分类 $\mathrm{EXT}(F_1,F_2)$ 的同构类. 若 $X$ 定义在域 $\Bbbk$ 上, 则从仿射 $\Bbbk$-概形范畴出发的函子 $T\mapsto\operatorname{Ext}_{\mathcal O_{X\times T}}^1(p_1^*F_1,p_1^*F_2)$ 由 $\mathbb A(\operatorname{Ext}_{\mathcal O_X}^1(F_1,F_2)^\vee)$ 表示 (习题~\bookref{1.3.22}). 群 $\underline{\operatorname{Aut}}(F_1)\times\underline{\operatorname{Aut}}(F_2)$ 在 $\mathbb A(\operatorname{Ext}_{\mathcal O_X}^1(F_1,F_2)^\vee)$ 上的作用由下式给出:
\[
\begin{aligned}
&(\beta_1,\beta_2)\cdot[0\to F_2\xrightarrow aE\xrightarrow bF_1\to0]\\
&\qquad=[0\to F_2\xrightarrow{a\circ\beta_2^{-1}}E\xrightarrow{\beta_1\circ b}F_1\to0].
\end{aligned}
\]
对角子群 $\mathbb G_m=\{(t,t)\}\subseteq\underline{\operatorname{Aut}}(F_1)\times\underline{\operatorname{Aut}}(F_2)$ 对应于 $E$ 上的数乘自同构, 在 $\mathbb A(\operatorname{Ext}_{\mathcal O_X}^1(F_1,F_2)^\vee)$ 上平凡作用. 另一方面, $\mathbb G_m=\{(t,1)\}\subseteq\underline{\operatorname{Aut}}(F_1)\times\underline{\operatorname{Aut}}(F_2)$ 通过数乘作用非平凡地作用于 $\mathbb A(\operatorname{Ext}_{\mathcal O_X}^1(F_1,F_2)^\vee)$. 对后一作用, $X\times\mathbb A(\operatorname{Ext}_{\mathcal O_X}^1(F_1,F_2)^\vee)$ 上的万有扩张是 $\mathbb G_m$-等变的, 并定义态射
\begin{equation}\label{src:9.2.11}
[\mathbb A(\operatorname{Ext}_{\mathcal O_X}^1(F_1,F_2)^\vee)/\mathbb G_m]\to\cBun(X),
\end{equation}
它在非分裂扩张的轨迹上限制为态射 $\mathbb P(\operatorname{Ext}_{\mathcal O_X}^1(F_1,F_2)^\vee)\to\cBun(X)$.

在相对情形下, 有:
\begin{exercise}[中等]\label{src:9.2.12}
设 $f:X\to S$ 是 Noether 概形之间的固有平坦态射, 例如对曲线 $C$ 取 $X=C\times S\to S$. 设 $F_1,F_2$ 是 $X$ 上向量丛, 对所有 $s\in S$, 有 $\Hom_{\mathcal O_{X_s}}(F_{1,s},F_{2,s})=0$.
\begin{enumerate}[(a)]
\item 证明 $\mathbb E:=(f_*\mathcal E\!xt_{\mathcal O_X}^1(F_1,F_2))^\vee$ 是 $S$ 上向量丛, 其构造与 $S$ 上的基变换交换.
\item 证明函子
\[
\begin{aligned}
\Sch/S&\to\Sets\\
(T\to S)&\mapsto\operatorname{Ext}_{\mathcal O_{X_T}}^1(F_{1,T},F_{2,T})
\end{aligned}
\]
由 $\mathbb A(\mathbb E)\to S$ 表示.

\textit{提示:} 利用等同 $\operatorname{Ext}_{\mathcal O_{X_T}}^1(F_{1,T},F_{2,T})=\mathrm H^0(X_T,\mathcal E\!xt_{\mathcal O_{X_T}}^1(F_{1,T},F_{2,T}))$; 它可由局部到整体谱序列
\[
\mathrm H^p(T,f_{T,*}\mathcal E\!xt_{\mathcal O_{X_T}}^q(F_{1,T},F_{2,T}))
\Rightarrow\operatorname{Ext}_{\mathcal O_{X_T}}^{p+q}(F_{1,T},F_{2,T})
\]
的低次项给出的正合列建立.
\item 证明 $\mathbb P(\mathbb E)\to S$ 表示函子 $\Sch/S\to\Sets$, 它将 $T\to S$ 映到对 $(\mathcal L,\mathcal E)$ 的同构类, 其中 $\mathcal L\in\mathrm{Pic}(T)$ 且 $\mathcal E\in\operatorname{Ext}_{\mathcal O_{X_T}}(F_{1,T},F_{2,T}\otimes\mathcal L_{X_T})$, 对所有 $t\in T$, $\mathcal E_t\in\operatorname{Ext}_{\mathcal O_{X_t}}(F_{1,t},F_{2,t})$ 非分裂. 这里 $(\mathcal L,\mathcal E)\sim(\mathcal L',\mathcal E')$ 表示存在同构 $\alpha:\mathcal L\xrightarrow{\sim}\mathcal L'$ 和 $\beta:\mathcal E\xrightarrow{\sim}\mathcal E'$, 使下图交换:
\[
\begin{tikzcd}[ampersand replacement=\&,column sep=small]
0\arrow[r]\&F_{2,T}\otimes\mathcal L_{X_T}\arrow[r]\arrow[d,"{\mathrm{id}\otimes\alpha_{X_T}}"']\&\mathcal E\arrow[r]\arrow[d,"\beta"]\&F_{1,T}\arrow[r]\arrow[d,"\mathrm{id}"]\&0\\
0\arrow[r]\&F_{2,T}\otimes\mathcal L'_{X_T}\arrow[r]\&\mathcal E'\arrow[r]\&F_{1,T}\arrow[r]\&0.
\end{tikzcd}
\]
此外, 证明在 $X\times_S\mathbb P(\mathbb E)$ 上存在万有扩张
\[
0\to p_1^*F_2\otimes p_2^*\mathcal O_{\mathbb P(\mathbb E)}(1)
\to\mathcal E_{\mathrm{univ}}\to p_1^*F_2\to0.
\]
\end{enumerate}
另见 \cite[附录 II, 命题 2]{Ses82} 和 \cite{Lan83a}.
\end{exercise}
\textbf{扩张叠.} 由于扩张有自同构, 将扩张的模对象描述为叠更为自然. 这使我们方便地处理 $f_*\mathcal E\!xt_{\mathcal O_X}^1(F_1,F_2)$ 不是向量丛的情形, 例如 $\operatorname{Ext}_{\mathcal O_X}^1(F_{1,s},F_{2,s})$ 的维数依赖于点 $s\in S$ 时.

在 $(\Sch/S)_{\text{ét}}$ 上定义叠
\[
\underline{\mathrm{EXT}}(F_1,F_2),
\]
其在 $T\to S$ 上的范畴为 $F_1,F_2$ 在 $X_T$ 上的基变换 $F_{1,T},F_{2,T}$ 的扩张群胚 $\mathrm{EXT}(F_{1,T},F_{2,T})$.
\begin{proposition}\label{src:9.2.13}
若 $f:X\to S$ 是曲线族, $F_1,F_2$ 是 $X$ 上的向量丛, 则 $\underline{\mathrm{EXT}}(F_1,F_2)$ 是代数叠, 同构于例~\bookref{4.9.33} 中定义的 $S$ 上的向量丛叠 $\mathfrak C(\mathrm Rf_*\mathcal H\!om_{\mathcal O_X}(F_1,F_2))$. 特别地, $\underline{\mathrm{EXT}}(F_1,F_2)\to S$ 光滑且拟紧, 其相对维数对任意点 $s\in S$ 计算为
\[
\dim\operatorname{Ext}_{\mathcal O_{X_s}}^1(F_{1,s},F_{2,s})
-\dim\operatorname{Ext}_{\mathcal O_{X_s}}^0(F_{1,s},F_{2,s}).
\]
\end{proposition}
\begin{proof}
设 $K^\bullet:K^0\xrightarrow dK^1$ 是计算 $\mathrm Rf_*\mathcal H\!om_{\mathcal O_X}(F_1,F_2)$ 的向量丛复形. 则
\[
\mathfrak C(\mathrm Rf_*\mathcal H\!om_{\mathcal O_X}(F_1,F_2))
\cong\mathfrak C(K^\bullet)\cong[\mathbb A(K^{1,\vee})/\mathbb A(K^{0,\vee})],
\]
其中加法群 $\mathbb A(K^{0,\vee})$ 通过 $v_0\cdot v_1=v_1+d(v_0)$ 在 $\mathbb A(K^{1,\vee})$ 上作用. 若 $g:T\to S$ 是来自仿射概形的态射, 则所有主 $\mathbb A(K^{0,\vee})$-丛平凡, 所以
\[
\begin{aligned}
\mathfrak C(K^\bullet)(T)&\cong[\mathbb A(K^{1,\vee})(T)/\mathbb A(K^{0,\vee})(T)]\\
&\cong[\mathrm H^0(T,g^*K^1)/\mathrm H^0(T,g^*K^0)]&&\text{(见 \bookeqref{4.9.34})}\\
&\cong\mathrm{EXT}(F_{1,T},F_{2,T}),
\end{aligned}
\]
最后一步来自 $\operatorname{Ext}_{\mathcal O_{X_T}}^i(F_{1,T},F_{2,T})\cong\mathrm H^i(g^*K^\bullet)$.
\end{proof}
\begin{remark}\label{src:9.2.14}
若每个 $s\in S$ 都满足 $\Hom(F_{1,s},F_{2,s})=0$, 则由上同调与基变换 (\bookref{A.6.8}), $f_*\mathcal E\!xt_{\mathcal O_X}^1(F_1,F_2)$ 是 $S$ 上的向量丛. 此时得到习题~\bookref{9.2.12}: $\mathrm{EXT}(F_1,F_2)$ 在 $S$ 上由 $\mathbb A(f_*\mathcal E\!xt_{\mathcal O_X}^1(F_1,F_2)^\vee)$ 表示. 更一般地, 若 $f_*\mathcal E\!xt_{\mathcal O_X}^i(F_1,F_2)$ 对 $i=0,1$ 都是向量丛, 则 $\underline{\mathrm{EXT}}(F_1,F_2)\to S$ 是带 $f_*\mathcal H\!om_{\mathcal O_X}(F_1,F_2)$-带的 gerbe.
\end{remark}
\begin{remark}[万有扩张]\label{src:9.2.15}
对两对秩次数 $(r_1,d_1)$ 和 $(r_2,d_2)$, 将此构造用于态射
\[
C\times\cBun_{r_1,d_1}\times\cBun_{r_2,d_2}\to\cBun_{r_1,d_1}\times\cBun_{r_2,d_2}
\]
及向量丛 $p_{12}^*\mathcal F_{1,\mathrm{univ}}$ 和 $p_{13}^*\mathcal F_{2,\mathrm{univ}}$, 其中 $\mathcal F_{i,\mathrm{univ}}$ 是 $C\times\cBun_{r_i,d_i}$ 上的万有向量丛. 得到扩张叠
\[
q:\underline{\mathrm{EXT}}(p_{12}^*\mathcal F_{1,\mathrm{univ}},p_{13}^*\mathcal F_{2,\mathrm{univ}})
\to\cBun_{r_1,d_1}\times\cBun_{r_2,d_2},
\]
其在 $T$ 上的对象为 $C_T$ 上的向量丛 $F_1,F_2$, 满足 $\operatorname{rk}F_i=r_i$, $\deg F_i=d_i$, 以及扩张 $0\to F_2\to E\to F_1\to0$; 位于 $T$ 的恒等映射之上的态射为如下交换图的数据:
\[
\begin{tikzcd}[ampersand replacement=\&,column sep=small]
0\arrow[r]\&F_2\arrow[r,"a"]\arrow[d,"\beta_2"']\&E\arrow[r,"b"]\arrow[d,"\alpha"]\&F_1\arrow[r]\arrow[d,"\beta_1"]\&0\\
0\arrow[r]\&F_2'\arrow[r,"a'"]\&E'\arrow[r,"b'"]\&F_1'\arrow[r]\&0,
\end{tikzcd}
\]
其中 $\alpha,\beta_1,\beta_2$ 都是同构.

在 $C\times\underline{\mathrm{EXT}}(p_{12}^*\mathcal F_{1,\mathrm{univ}},p_{13}^*\mathcal F_{2,\mathrm{univ}})$ 上存在秩 $r=r_1+r_2$, 次数 $d=d_1+d_2$ 的万有扩张 $\mathcal E_{\mathrm{univ}}$:
\[
0\to(\mathrm{id}\times q)^*p_{13}^*\mathcal F_{2,\mathrm{univ}}
\to\mathcal E_{\mathrm{univ}}\to(\mathrm{id}\times q)^*p_{12}^*\mathcal F_{1,\mathrm{univ}}\to0.
\]
这定义一个可表态射 $\underline{\mathrm{EXT}}(p_{12}^*\mathcal F_{1,\mathrm{univ}},p_{13}^*\mathcal F_{2,\mathrm{univ}})\to\cBun_{r,d}$, 嵌入下图:
\[
\begin{tikzcd}[ampersand replacement=\&,column sep=small]
\underline{\mathrm{EXT}}(p_{12}^*\mathcal F_{1,\mathrm{univ}},p_{13}^*\mathcal F_{2,\mathrm{univ}})\arrow[r,"{[\mathcal E_{\mathrm{univ}}]}"]\arrow[d,"q"']\&\cBun_{r,d}\\
\cBun_{r_1,d_1}\times\cBun_{r_2,d_2}\&{}
\end{tikzcd}
\qquad
\begin{tikzcd}[ampersand replacement=\&,column sep=small]
{[0\to F_2\to E\to F_1\to0]}\arrow[r,mapsto]\arrow[d,mapsto]\&{[E]}\\
{([F_1],[F_2])}\&{}
\end{tikzcd}
\]
\end{remark}
